\subsection{The Cauchy criterion}

The definition of convergence refers to a candidate limit \(L\). Sometimes we want to decide whether a sequence converges before knowing that limit. The Cauchy criterion compares late terms directly with one another.

\begin{definitionbox}[Cauchy sequence]
A sequence \((a_n)\) is a Cauchy sequence if, for every \(\varepsilon>0\), there exists \(N\in\mathbb N\) such that
\[
m,n\ge N \quad\Longrightarrow\quad |a_m-a_n|<\varepsilon.
\]
\end{definitionbox}

Thus, sufficiently late terms must be arbitrarily close to each other. The definition does not name the value they approach.

\begin{theorembox}[Cauchy criterion in \(\mathbb R\)]
A real sequence converges if and only if it is a Cauchy sequence.
\end{theorembox}

\begin{proofbox}[Why convergence implies the Cauchy property]
Suppose \(a_n\to L\). Given \(\varepsilon>0\), choose \(N\) so that
\[
n\ge N \quad\Longrightarrow\quad |a_n-L|<\frac\varepsilon2.
\]
Then for \(m,n\ge N\), the triangle inequality gives
\[
|a_m-a_n|
\le |a_m-L|+|a_n-L|
<\frac\varepsilon2+\frac\varepsilon2
=\varepsilon.
\]
Hence every convergent sequence is Cauchy.
\end{proofbox}

The reverse implication uses the completeness of the real numbers: real Cauchy sequences cannot approach a missing point outside \(\mathbb R\).

\begin{examplebox}[A sequence that is not Cauchy]
For \(a_n=(-1)^n\), choose arbitrarily large even \(m\) and odd \(n\). Then
\[
|a_m-a_n|=|1-(-1)|=2.
\]
The late terms do not become close to one another, so the sequence is not Cauchy and therefore does not converge.
\end{examplebox}

\begin{interpretationbox}[Internal test for convergence]
The ordinary definition compares terms with an external target \(L\). The Cauchy criterion tests whether the tail has internally settled down enough that a limit must exist.
\end{interpretationbox}